\documentclass[../../../main.tex]{subfiles}

\begin{document}

\chapter{C Programming Language}

I always wanted to learn C. I mean C really is the fundamental
language that inspired so many others like PHP, Java, C++, and
so much more. Also there is that coolness in C. Is it only me?
Or do you also feel that sensation when you see 90+\% gray dot
in GitHub repository for Linux Kernel? Writing in C and going
low level has its own coolness. In this note, we will discuss some
fundamental features in C for us to actually use it for writing
goodies. Many of the contents here are based on the video
``C Programming Tutorial for Beginners'' from freeCodeCamp by
Mike Dane and many other online tutorials \cite{CS03-1}.

\section{Foundations}

The previous note contains much more detailed explanations on
data structures, algorithms, and solving problems. So in this
section, instead of discussing everything from scratch again,
let's focus on the syntax for C and some distinct features that
C shines with.

For our first program, let's output ``Hello, world!'' and see how
to compile C programs.

\begin{cbox}{26clcwy-01}{26clcwy-01}
\begin{lstlisting}[style=cstyle]
#include <stdio.h>

int main() {
    printf("Hello, world!\n");

    return 0;
}
\end{lstlisting}
\end{cbox}

For our header, we have standard input/output and we use \verb|printf|
to print output ``Hello, world!''. Now there are two popular ways of
compiling it and we can use either GCC or Clang. They are nearly
identical (please don't get mad at me!), and I will be using Clang
since I am on macOS.

\begin{basicbox}
\begin{lstlisting}[style=basic]
$ clang -Wall program.c -o program
$ ./program
Hello, world!
\end{lstlisting}
\end{basicbox}

With this, let's discuss about input/output and basic data types.

\subsection{Input/Output and Basic Data Types}

Now this \verb|printf| function can be very powerful. Like Java,
we can pass values into placeholders like \%d, \%f, \%s, \%c,
namely double, float, string, and character. Here is an example.

\begin{cbox}{26clcwy-02}{26clcwy-02}
\begin{lstlisting}[style=cstyle]
#include <stdio.h>

int main() {
    int age = 123;
    double gpa = 3.14;
    char name[] = "Enoch";

    printf("%s is %d years old.\n", name, age);
    printf("He has a GPA of %f and a letter grade of %c\n", gpa, 'A');
    printf("%s also loves math. For example...\n", name);
    printf("%f / %d = %f\n", 1.0, 5, 1.0 / 5);

    return 0;
}
\end{lstlisting}
\end{cbox}

As you can see, we can define integer, double, float, and character
as usual. We can also pass our values into the place holder.
You may have noticed that we wrote \verb|char name[] = "Enoch";|
instead of \verb|string name = "Enoch";|. This is because unlike
other languages, C doesn't have data type string in its usual sense.
We will discuss more about this as we discuss arrays and pointers.
Finally, notice that we wrote $1.0$ instead of $1$ for our first
\verb|%f| in the last line. This is because $1$ by itself is an
integer and we will get an error. Compiling and running, we get the
following output.

\begin{outputbox}
\begin{lstlisting}[style=basic]
Enoch is 123 years old.
He has a GPA of 3.140000 and a letter grade of A
Enoch also loves math. For example...
1.000000 / 5 = 0.200000
\end{lstlisting}
\end{outputbox}

As you can see, we get $3.140000$ instead of $3.14$. When displaying,
C likes to give a more accurate value. Continuing, we have same
commenting style as C++.

\begin{cbox}{26clcwy-03}{26clcwy-03}
\begin{lstlisting}[style=cstyle]
#include <stdio.h>

int main() {
    /*
     * This is a multi-line comment
     * This is fun!
     */

    printf("Hi there\n");
    // printf("This is an inline comment");
    
    return 0;
}
\end{lstlisting}
\end{cbox}
\begin{outputbox}
\begin{lstlisting}[style=basic]
Hi there
\end{lstlisting}
\end{outputbox}

Similar to other languages, we can define a constant with \verb|const|
and arrays with \verb|arr[]| as well. Here is an example of using
constants and arrays.

\begin{cbox}{26clcwy-04}{26clcwy-04}
\begin{lstlisting}[style=cstyle]
#include <stdio.h>

int main() {
    const int FIXED_VALUE = 5;
    int arr[5] = { 1, 2, 3, 4, 5 };
    double tensor[2][3] = {{0.0, 0.1, 0.2},
                           {1.0, 1.1, 1.2}};
    char string[10] = "string";

    printf("Fixed value: %d\n", FIXED_VALUE);
    printf("Second element of the array: %d\n", arr[1]);
    printf("Second row, third column: %f\n", tensor[1][2]);
    printf("Third charater in the string: %c\n", string[2]);

    return 0;
}
\end{lstlisting}
\end{cbox}
\begin{outputbox}
\begin{lstlisting}[style=basic]
Fixed value: 5
Second element of the array: 2
Second row, third column: 1.200000
Third charater in the string: r
\end{lstlisting}
\end{outputbox}

Now as mentioned previously, strings are quite different in C.
They are not totally different data types, but just an array of
characters. This simple nature of C is one of the factors that makes
C special, but also slightly annoying.

Now before moving onto inputs and outputs, here is one more data type.
This one called \textbf{struct} is quite interesting because it is
user-defined and also a composite data type. It can look like a class
without functions, but they are different. Struct defines data type
that we can use to organize information. Here is an example usage.

\begin{cbox}{26clcwy-05}{26clcwy-05}
\begin{lstlisting}[style=cstyle]
#include <stdio.h>

typedef struct {
    float a;
    float b;
    float c;
} array ;

struct vector {
    float x;
    float y;
};

int main() {
    // Using vec
    struct vector vec;
    vec.x = 0.1;
    vec.y = 0.2;
    printf("v   = <%f, %f>\n", vec.x, vec.y);

    // Using arr
    array arr = { .a = 0.1, .b = 0.2, .c = 0.3, };
    printf("arr = [%f, %f, %f]\n", arr.a, arr.b, arr.c);

    return 0;
}
\end{lstlisting}
\end{cbox}
\begin{outputbox}
\begin{lstlisting}[style=basic]
v   = <0.100000, 0.200000>
arr = [0.100000, 0.200000, 0.300000]
\end{lstlisting}
\end{outputbox}

Struct is nice because we can use it like a template and reuse it
multiple times. Also we will discuss more on \verb|typedef| later
but we can use it as shown above to avoid writing more characters.
Let's now turn our attention to inputs and outputs in C. We already
discussed outputs, so let's see how we can take values in from the
user in C.

\begin{cbox}{26clcwy-06}{26clcwy-06}
\begin{lstlisting}[style=cstyle]
#include <stdio.h>

int main() {
    char name[10], fullName[30], alp;
    int age;
    double date;

    printf("Full Name: ");
    fgets(fullName, 30, stdin);
    printf("First Name: ");
    scanf("%s", name); // No ampersand here
    printf("Age: ");
    scanf("%d", &age); // You need this ampersand!
    printf("Date: ");
    scanf("%lf", &date); // %lf not %f
    printf("Alphabet: ");
    scanf(" %c", &alp); // You need the space in front of %c!

    printf("Hi %s, today is %f\n", name, date);
    printf("Your full name is %s", fullName);
    printf("You are %d years old and you chose %c!\n", age, alp);

    return 0;
}
\end{lstlisting}
\end{cbox}
\begin{outputbox}
\begin{lstlisting}[style=basic]
Full Name: John Doe
First Name: John
Age: 100
Date: 11.18
Alphabet: J
Hi John, today is 11.180000
Your full name is John Doe
You are 100 years old and you chose J!
\end{lstlisting}
\end{outputbox}

The ampersand that you see is there because we pass a pointer to
\verb|scanf|, but we will get there in a second. Also when passing
a character, float, and a string with spaces in it, we need to
consider few factors. This is the place where I said its kinda
annoying and awkward at first, but they are what makes C, the C!

Also, C is designed to give us low level access, so as you can see
with getting string inputs, we can manually allocate memory necessary.
Now continuing with input/output and simple data types, let's move on
to conditionals and loops.

\subsection{Conditionals and Loops}

The syntax for conditionals for C is identical to the ones for C++.
Here is an example of using if-else statements.

\begin{cbox}{26clcwy-07}{26clcwy-07}
\begin{lstlisting}[style=cstyle]
#include <stdio.h>

int main() {
    if (1 > 2) {
        printf("1 > 2\n");
    } else if (3 > 2) {
        printf("3 > 2\n");
    }

    if (1 > 2 || 3 > 2) {
        printf("1 > 2 or 3 > 2\n");
    }
    if (1 < 2 && 3 > 2) {
        printf("1 < 2 and 3 > 2\n");
    }

    return 0;
}
\end{lstlisting}
\end{cbox}
\begin{outputbox}
\begin{lstlisting}[style=basic]
3 > 2
1 > 2 or 3 > 2
1 < 2 and 3 > 2
\end{lstlisting}
\end{outputbox}

Similarly, we could also use switch statements.

\begin{cbox}{26clcwy-08}{26clcwy-08}
\begin{lstlisting}[style=cstyle]
#include <stdio.h>

int main() {
    int choice;
    printf("Your choice (1-3): ");
    scanf("%d", &choice);

    switch (choice) {
        case 1:
            printf("Your choice is 1\n");
            break;
        case 2:
            printf("Your choice is 2\n");
            break;
        case 3:
            printf("Your choice is 3\n");
            break;
        default:
            printf("Invalid Option\n");
            break;
    }

    return 0;
}
\end{lstlisting}
\end{cbox}
\begin{outputbox}
\begin{lstlisting}[style=basic]
$ clang -Wall program.c -o program
$ ./program
Your choice (1-3): 2
Your choice is 2
$ clang -Wall program.c -o program
$ ./program
Your choice (1-3): 4
Invalid Option
\end{lstlisting}
\end{outputbox}

Continuing, we also have the same syntax for for, while, and do-while
loops.

\begin{cbox}{26clcwy-09}{26clcwy-09}
\begin{lstlisting}[style=cstyle]
#include <stdio.h>

int main() {
    for (int i = 0; i < 5; i++) {
        printf("%d ", i);
    }
    printf("\n");

    int ctr = 0;
    while (ctr < 5) {
        printf("%d ", ctr);
        ctr++;
    }
    printf("\n");

    do {
        printf("%d ", ctr);
        ctr++;
    } while (ctr < 0);
    printf("\n");

    return 0;
}
\end{lstlisting}
\end{cbox}
\begin{outputbox}
\begin{lstlisting}[style=basic]
0 1 2 3 4
0 1 2 3 4
5
\end{lstlisting}
\end{outputbox}

Now before wrapping up this section, let's discuss about functions.

\subsection{Functions}

Functions in C works like many other languages. Here is an example
of a hello user function.

\begin{cbox}{26clcwy-10}{26clcwy-10}
\begin{lstlisting}[style=cstyle]
#include <stdio.h>

void helloUser(char name[]) {
    printf("Hello %s\n", name);
}

int main() {
    char name[20];
    printf("Your name: ");
    scanf("%s", name);
    helloUser(name);

    return 0;
}
\end{lstlisting}
\end{cbox}
\begin{outputbox}
\begin{lstlisting}[style=basic]
$ clang -Wall program.c -o program
$ ./program
Your name: Enoch
Hello Enoch
\end{lstlisting}
\end{outputbox}

When we want our main function to be on the very top, then we need
something called prototype. Function prototypes and declaration are
used interchangeably with slight difference in nuance. As shown below,
we need to basically tell the compiler that the function that we will
be using in the main function exists, so that it doesn't complain
about it.

\begin{cbox}{26clcwy-11}{26clcwy-11}
\begin{lstlisting}[style=cstyle]
#include <stdio.h>

double square(double num);

int main () {
    printf("1.4^2 is equal to %f\n", square(1.4));

    return 0;
}

double square(double num) {
    return num * num;
}
\end{lstlisting}
\end{cbox}
\begin{outputbox}
\begin{lstlisting}[style=basic]
1.4^2 is equal to 1.960000
\end{lstlisting}
\end{outputbox}

This is it for this section! In the next section, we will discuss
more about pointers and we mentioned briefly here.

\end{document}


